\section{Synthesis and what carries forward}\label{m:sec:summary}

The chapter has assembled three layers of machinery.  The first is the standard
Markov-chain apparatus: exponential clocks and competing events, continuous-time
generators, Kolmogorov equations, probability generating functions and the linear
birth--death process.  Its generating-function equation
\cref{m:eq:bdpde} is the template for the later calculations.

The second layer is the branching and quasi-stationary picture.  The binary
Galton--Watson process supplies an exact survival recursion, its critical scaling
and its total-progeny law.  In the subcritical regime, conditioning on survival
produces a limiting mean \(1/A(p)\) and variance \cref{m:eq:condvar}.  Continuous
time provides the explicit comparison constant
\(A_{\mathrm c}(p)=(1-2p)/(1-p)\).  The discrete constant is used here as a named
object only: its product, series, bounds, near-critical asymptotics,
Koenigs-function identification and obstruction theory are all reserved for
\ChA.

The third layer is the method of characteristics for bivariate generating-function
equations.  It was checked against absorption--death, where an independent-particle
answer was already available, and then applied to absorption--birth--death to give
the closed form \cref{m:eq:Gabd}.  The appendices show how the special-function
identity enters and how state probabilities are recovered from the resulting PGFs.

Two questions remain naturally on the mathematical side of this chapter.  First,
when extinction and catastrophe coexist, the relevant conditional law is that of a
killed chain conditioned on avoiding both mechanisms; the eigenrelation in
\cref{m:eq:bdc-qsd-eigenrelation} gives the correct starting point, but not a
general solution.  Second, the Jensen lower bound for random catastrophe exposure
is rigorous, while its tightness depends on fluctuations of the accumulated
population and must be assessed model by model.  Later modelling chapters take up
the birth--death--catastrophe and multi-compartment extensions; this chapter stops
with the reusable toolkit.
